% !TeX root = ./settheory-screen.tex
% settheory.tex
%
% driver file settheory.tex to produce text

\pretocmd{\OLEndChapterHook}{%
  \IfFileExists{include/summary-\thechapter}{%
    \section*{Resumo}\addcontentsline{toc}{section}{Resumo}%
    \let\emph\textbf\input{include/summary-\thechapter}\let\emph\textit%
  }{}%
}{}{}

\problemsperchapter

\allowdisplaybreaks

\frontmatter

\OLPfrontmatter

\chapter{Sobre este Livro}

Este livro é um Recurso Educacional Aberto. Ele foi escrito para estudantes
com algum conhecimento de lógica e de matemática do ensino médio.
Seu objetivo é oferecer uma primeira visão da filosofia da teoria dos
conjuntos. Ao final deste livro, os estudantes que o lerem podem ter uma
noção:
\begin{enumerate}
	\item de por que a teoria dos conjuntos surgiu;
	\item de como incorporar grandes partes da matemática em teoria dos conjuntos + aritmética;
	\item de como incorporar a própria aritmética na teoria dos conjuntos;
	\item do que envolve a concepção cumulativa iterativa de conjunto;
	\item de como se pode tentar justificar os axiomas de ZFC.
\end{enumerate}
O livro surgiu de um curso curto que eu ministrei no Departamento de Filosofia de Cambridge. Antes de mim, ele foi ministrado por Luca Incurvati e Michael
Potter. Ao escrever este livro---e, de forma mais geral, o curso---fiquei
profundamente em dívida com Luca e Michael. Espero que isso fique claro no próprio texto do livro; mas também quero registrar aqui
meus sinceros agradecimentos a ambos.

A maior parte deste livro foi originalmente publicada como \emph{Open
Set Theory}; este livro é seu sucessor. Contribuí com todo o material deste livro para o \href{http://openlogicproject.org}{\emph{Open Logic Project}}, onde ele está disponível gratuitamente. (\Crefrange{sfr:set::chap}{sfr:siz::chap} foram elaborados, com apenas
pequenas alterações, a partir de material já existente no OLP; essas alterações agora também fazem parte do OLP.) Veja
\href{http://openlogicproject.org/}{openlogicproject.org} para mais
informações.

\begin{flushright}
Tim Button
\\University College London
\\Outubro de 2021
\end{flushright}


\mainmatter

\part{Prelúdio}

\olchapter{his}{set}{História e Mitologia}

Para compreender o significado filosófico da teoria dos conjuntos, ajuda
ter alguma ideia de por que ela surgiu. Para isso, ajuda pensar um pouco
na história e na mitologia da matemática. Assim, antes mesmo de
começarmos a discutir a teoria dos conjuntos, começaremos com uma
``história'' muito breve. Mas a colocamos entre aspas, porque é muito breve,
extremamente seletiva e algo contestável.

\olimport*[history/set-theory]{infinitesimals}
\olimport*[history/set-theory]{limits}
\olimport*[history/set-theory]{pathologies}
\olimport*[history/set-theory]{mythology}

\section{Roteiro}

Parte da lição da secção anterior é que a história da matemática foi em grande
medida escrita pelos vencedores. Tinham machados para afiar; machados
filosóficos e matemáticos. O estudo sério da história da matemática é
seriamente difícil (e gratificante), e a Coruja de Minerva só levanta
voo ao crepúsculo.

Apesar disso, é incontestável que os resultados ``patológicos''
envolveram o desenvolvimento de novas ferramentas matemáticas fascinantes
e uma reconsideração dos padrões de rigor matemático. Por exemplo,
exigiram pensar o contínuo (a ``reta real'') de determinado modo,
e pensar as funções como mapas ponto a ponto. E, no fim, o
desenvolvimento completo de todas essas ferramentas exigiu o
desenvolvimento rigoroso da teoria dos conjuntos. O restante deste livro
explicará algo desse desenvolvimento.

\Olref[sfr][][]{part} apresentará uma versão da teoria \emph{ingênua} dos
conjuntos, facilmente suficiente para desenvolver toda a matemática
acabada de descrever. Isto levará algum tempo. Mas, ao fim de
\olref[sfr][][]{part}, estaremos em condições de compreender como tratar
números reais como certos conjuntos, e como tratar funções sobre
eles---incluindo curvas preenchedoras de espaço---como \emph{outros}
conjuntos.

Mas a \emph{ingenuidade} dessa teoria dos conjuntos emergirá na
\olref[sth][][]{part}, quando encontrarmos paradoxos da teoria dos conjuntos
e a necessidade sentida de descrever as coisas com muito mais precisão.
Nesse ponto, precisaremos desenvolver um tratamento axiomático dos
conjuntos, com o qual possamos recuperar todos os nossos resultados
ingênuos, evitando (esperemos) paradoxos. (A Coruja do rigor matemático
também só levanta voo ao crepúsculo.)

\OLEndChapterHook

\olpart{sfr}{Teoria Ingênua dos Conjuntos}

\section*{Introdução à \Olref[sfr][][]{part}}

Na \olref[sfr][][]{part}, consideraremos os conjuntos de forma ingênua e
informal. \Olref[sfr][set][]{chap} apresentará a ideia básica, a de que os
conjuntos são coleções consideradas extensionalmente, e introduzirá algumas
operações muito elementares. Em seguida,
\crefrange{sfr:rel::chap}{sfr:fun::chap} explicarão como a teoria dos conjuntos
nos permite falar de relações e, portanto, de funções.

\Crefrange{sfr:siz::chap}{sfr:infinite::chap} tratarão então de algumas das
primeiras realizações da teoria ingênua dos conjuntos. Em
\olref[sfr][siz][]{chap}, exploramos como comparar conjuntos quanto ao seu
tamanho. No \olref[sfr][arith][]{chap}, exploramos como se podem reduzir os
inteiros, os racionais e os reais à teoria dos conjuntos juntamente com
aritmética elementar. No \olref[sfr][infinite][]{chap}, consideramos como se
pode implementar aritmética elementar no seio da teoria dos conjuntos.

Repetindo, tudo isto será feito de forma \emph{ingênua}. Mas tudo o que
fizermos na \olref[sfr][][]{part} pode ser feito com rigor perfeito, na teoria
formal dos conjuntos que introduzimos na \olref[sth][][]{part}.

\olchapter{sfr}{set}{Começando}

\olimport*[sets-functions-relations/sets]{basics}
\olimport*[sets-functions-relations/sets]{subsets}
\olimport*[sets-functions-relations/sets]{important-sets}
\olimport*[sets-functions-relations/sets]{unions-and-intersections}
\olimport*[sets-functions-relations/sets]{pairs-and-products}
\olimport*[sets-functions-relations/sets]{russells-paradox}

\OLEndChapterHook

\olchapter{sfr}{rel}{Relações}

\olimport*[sets-functions-relations/relations]{relations-as-sets}
\olimport*[sets-functions-relations/relations]{reflections}
\olimport*[sets-functions-relations/relations]{special-properties}
\olimport*[sets-functions-relations/relations]{equivalence-relations}
\olimport*[sets-functions-relations/relations]{orders}
\olimport*[sets-functions-relations/relations]{operations}

\OLEndChapterHook

\olchapter{sfr}{fun}{Funções}

\olimport*[sets-functions-relations/functions]{function-basics}
\olimport*[sets-functions-relations/functions]{function-kinds}
\olimport*[sets-functions-relations/functions]{functions-relations}
\olimport*[sets-functions-relations/functions]{inverses}
\olimport*[sets-functions-relations/functions]{composition}

\OLEndChapterHook

\olchapter{sfr}{siz}{O Tamanho dos Conjuntos}

% Igualar espaço vertical dos displays ao do texto «longo»: em PT as quebras
% de linha mudam o \predisplaysize e o TeX usa skips «curtos» (menores) só
% nalguns displays — o cap. 5 fica visualmente irregular face ao original.
\begingroup
\setlength{\abovedisplayshortskip}{\abovedisplayskip}%
\setlength{\belowdisplayshortskip}{\belowdisplayskip}%

\olimport*[sets-functions-relations/size-of-sets]{introduction}
\olimport*[sets-functions-relations/size-of-sets]{enumerability-alt}
\olimport*[sets-functions-relations/size-of-sets]{zig-zag}
\olimport*[sets-functions-relations/size-of-sets]{pairing}
\olimport*[sets-functions-relations/size-of-sets]{non-enumerability-alt}
\olimport*[sets-functions-relations/size-of-sets]{reduction-alt}
\olimport*[sets-functions-relations/size-of-sets]{equinumerous-sets}
\olimport*[sets-functions-relations/size-of-sets]{comparing-size}
\olimport*[sets-functions-relations/size-of-sets]{schroder-bernstein}
\olimport*[history/set-theory]{cantor-plane}
\olimport*[history/set-theory]{hilbert-curve}

\endgroup

\OLEndChapterHook

\olsetchapter{nosection}
\chapter{Aritmetização}

No \olref[his][set][]{chap}, consideramos parte do enquadramento
histórico, nomeadamente \emph{por que} temos sequer teoria dos
conjuntos.
\Crefrange{sfr:set::chap}{sfr:siz::chap} percorreram depois alguns
princípios da teoria ingênua dos conjuntos. Compreendemos assim, de
modo ingênuo, como construir relações e funções, comparar os tamanhos
dos conjuntos e \emph{coisas desse gênero}.

Com isto assimilado, podemos abordar algumas das primeiras conquistas
da teoria dos conjuntos na redução (em certo sentido) de grandes
parcelas da matemática à teoria dos conjuntos e à aritmética. Esse é o
objetivo deste capítulo.

\olimport*[sets-functions-relations/arithmetization]{arithmetization}

\chapter{Conjuntos Infinitos}

No capítulo anterior, mostramos como construir, de certo modo, várias
coisas---inteiros, racionais e reais---assumindo uma teoria ingênua dos
conjuntos e os números naturais. A pergunta deste capítulo é: podemos
construir o próprio conjunto dos números naturais \emph{em si} usando a
teoria dos conjuntos?

\olimport*[sets-functions-relations/infinite]{infinite}
\olresetchapter

\OLEndPartHook

\olpart{sth}{A Concepção Iterativa}

\section*{Introdução à \Olref[sth][][]{part}}

\Olref[sfr][][]{part} tratou os conjuntos de modo ingênuo e informal. É
tempo de apertar o rigor e apresentar uma teoria formal dos conjuntos.
Esse é o objetivo da \olref[sth][][]{part}.

A nossa teoria formal é uma teoria de primeira ordem com um único
predicado binário, $\in$. Enunciaremos vários axiomas que regem o
comportamento da relação de pertença. Contudo, introduziremos estes
axiomas só à medida que forem necessários, e consideraremos como
poderíamos justificá-los quando os encontrarmos. Como resultado,
introduziremos os nossos axiomas \emph{ao longo} dos capítulos
seguintes.

Pode, no entanto, ser útil ao leitor ter uma lista de todos os axiomas
num só lugar. Eis \emph{todos} os axiomas que consideraremos na
\olref[sth][][]{part}. Como na \olref[sfr][][]{part}, a escolha de
minúsculas e maiúsculas nas letras orienta-se apenas pela legibilidade:

\begin{defish}
\emph{Extensionalidade.}
\[\forall A \forall B(\forall x(x \in A \liff x \in B) \lif A = B)\]
\end{defish}

\begin{defish}
\emph{União.}
\[
\forall A \exists U \forall x(x \in U \liff (\exists b \in A)x \in b),
\]
isto é, $\bigcup A$ existe para qualquer conjunto~$A$.
\end{defish}

\begin{defish}
\emph{Pares.}
\[
\forall a \forall b \exists P \forall x (x \in P \liff
(x = a \lor x = b)),
\]
isto é, $\{a, b\}$ existe para quaisquer $a$ e~$b$.
\end{defish}

\begin{defish}
\emph{Conjuntos potência.}
\[
\forall A \exists P \forall x(x \in P \liff (\forall z \in x)z \in A),
\]
isto é, $\Pow{A}$ existe para qualquer conjunto~$A$.
\end{defish}

\begin{defish}
\emph{Infinito.}
\begin{multline*}
\exists I((\exists o \in I)\forall x(x \notin o) \land {}\\
(\forall x \in I)(\exists s \in I)\forall z(z \in s \liff (z \in x \lor z = x))),
\end{multline*}
isto é, existe um conjunto que tem $\emptyset$ como !!a{element} e que
é fechado para $x \mapsto x \cup \{x\}$.
\end{defish}

\begin{defish}
\emph{Fundamentação.}
\[
\forall A(\forall x\, x \notin A \lor (\exists b \in A)(\forall x \in A)x \notin b),
\]
isto é, $A = \emptyset$ ou $(\exists b \in A)A\cap b = \emptyset$.
\end{defish}

\begin{defish}
\emph{Boa ordenação.}
Para todo o conjunto $A$, existe uma relação que bem ordena $A$.
(Escrever este axioma em lógica de primeira ordem seria demasiado penoso
para valer a pena aqui.)
\end{defish}

\begin{defish}
\emph{Esquema de separação.}
Para qualquer fórmula $\phi(x)$ que não contenha~``$S$'':
\[
\forall A \exists S \forall x(x \in S \liff (\phi(x) \land x \in A)),
\]
isto é, $\Setabs{x \in A}{\phi(x)}$ existe para qualquer conjunto~$A$.
\end{defish}

\begin{defish}
\emph{Esquema de substituição.} Para qualquer fórmula $\phi(x, y)$ que não contenha~``$B$'':
\[
\forall A((\forall x \in A)\lexists![y][\phi(x,y)] \lif \exists B \forall y (y \in B \liff (\exists x \in A)\phi(x,y))),
\]
isto é, $\Setabs{y}{(\exists x \in A)\phi(x,y)}$ existe para qualquer $A$, se $\phi$ for ``funcional''.
\end{defish}

Em ambos os esquemas, as fórmulas podem conter parâmetros. De fato,
ao longo de \olref[sth][][]{part}, seguimos a convenção de que qualquer
fórmula pode conter parâmetros. (Ver
\olref[sfr][infinite][induction]{sec} para recordar o que significa dizer
que uma fórmula pode conter parâmetros.)

\begin{description}
\item[$\Zminus$] é Extensionalidade, União, Pares, Conjuntos potência, Infinito, Separação.
\item[$\Z$] é $\Zminus$ mais Fundamentação.
\item[$\ZFminus$] é $\Z$ mais Substituição.
\item[$\ZF$] é $\ZFminus$ mais Fundamentação.
\item[$\ZFC$] é $\ZF$ mais Boa ordenação.
\end{description}

\olimport*[set-theory/story]{story}

\chapter{Passos rumo a $\Z$}

No capítulo anterior, consideramos a concepção iterativa de conjunto.
Neste capítulo, tentaremos extrair a maior parte dos axiomas da teoria de
conjuntos de Zermelo, isto é,~$\Z$. A abordagem é \emph{inteiramente}
inspirada em \citet{Boolos1971}, \citet{Scott1974} e
\citet{Shoenfield:AST}. 

\olsetchapter{nosection}
\olimport*[set-theory/z]{z}
\olresetchapter

\olimport*[set-theory/ordinals]{ordinals}

\olimport*[set-theory/spine]{spine}

\olimport*[set-theory/replacement]{replacement}

\olimport*[set-theory/ord-arithmetic]{ord-arithmetic}

\olimport*[set-theory/cardinals]{cardinals}

\chapter{Aritmética Cardinal}

No \olref[sth][cardinals][]{chap}, desenvolvemos uma teoria dos
cardinais. O passo seguinte é esboçar uma teoria da aritmética cardinal.
Este capítulo resume brevemente alguns fatos elementares e depois
aponta algumas das dificuldades, que se revelam fascinantes e
filosoficamente densas.

\olsetchapter{nosection}
\olimport*[set-theory/card-arithmetic]{card-arithmetic}
\olresetchapter

\olimport*[set-theory/choice]{choice}

\stopproblems
\def\ifproblems#1{}

\appendix

\setlength{\olphotowidth}{.45\textwidth}

\chapter{Biografias}

\olimport*[history/biographies]{georg-cantor}

\olimport*[history/biographies]{kurt-goedel}

\olimport*[history/biographies]{bertrand-russell}

\olimport*[history/biographies]{alfred-tarski}

\olimport*[history/biographies]{ernst-zermelo}

\backmatter

\clearpage

\photocredits

\bibliographystyle{\olpath/bib/natbib-oup}
\bibliography{\olpath/bib/open-logic}

% Lista de siglas e abreviaturas editoriais (curso Set Theory PT)

\clearpage
\phantomsection
\chapter*{Siglas e abreviaturas}
\markboth{\MakeUppercase{Siglas e abreviaturas}}{\MakeUppercase{Siglas e abreviaturas}}
\addcontentsline{toc}{chapter}{Siglas e abreviaturas}

\noindent
Abreviaturas editoriais e siglas de teorias listadas abaixo. Símbolos formais e convenções de escrita estão no capítulo \emph{Notação}; vocabulário matemático em prosa, no \emph{Glossário}.

\begin{description}
\item[Cf.] \emph{Conferir} (latim \textit{confer}): indica ao leitor que deve comparar ou consultar outra passagem.
\item[O.s.] \emph{Ou seja} (equivalente habitual do inglês \textit{i.e.} / \textit{id est}).
\item[p.ex.; P.ex.] \emph{Por exemplo} (equivalente do inglês \textit{e.g.}); \emph{P.ex.} reserva-se ao início de frase.
\item[ca.] \textit{Circa}: por volta de (datas, sobretudo em créditos de imagens).
\item[sse] \emph{Se e somente se} (equivalente do inglês \textit{iff} / \textit{if and only if}).
\item[$\Z$] Teoria de Zermelo \emph{com} o axioma de Fundação.
\item[$\Zminus$] Teoria de Zermelo \emph{sem} o axioma de Fundação.
\item[$\ZF$, $\ZFminus$] Teoria de Zermelo--Fraenkel; a variante $\ZFminus$ dispensa Fundação.
\item[$\ZFC$] Zermelo--Fraenkel com o Axioma da Escolha.
\item[$\LT$] Teoria dos Níveis (\textit{Level Theory}).
\item[$\Zr$] Extensão de $\LT$ com axiomas de Infinito e Conjunto potência, na qual todo o conjunto tem posto.
\item[$\Th{T}$] Letra ou nome curto para uma \emph{teoria} (conjunto de sentenças) de primeira ordem com símbolo primitivo~$\in$.
\item[HC] Hipótese do Contínuo.
\item[HGC] Hipótese Generalizada do Contínuo.
\item[(H)HC] Referência conjunta à HC ou à HGC quando o raciocínio é paralelo.
\item[TS] \emph{Tarski--Scott}: linha de definição de cardinalidade que evita o teorema do bom ordenamento (capítulo correspondente).
\end{description}

% Notação matemática (curso Set Theory PT) — alinhada ao texto importado em settheory.tex

\clearpage
\phantomsection
\chapter*{Notação}
\markboth{\MakeUppercase{Notação}}{\MakeUppercase{Notação}}
\addcontentsline{toc}{chapter}{Notação}

\noindent
Símbolos e abreviaturas usados com frequência neste volume (incluindo capítulos de pré-requisitos sobre conjuntos, relações e funções). Quando o mesmo símbolo tiver mais de um uso, o contexto (ou a definição na secção indicada) desambigua. Para vocabulário e noções em prosa, ver o \emph{Glossário}.

\section*{Lógica e teorias}

\begin{description}
\item[$\Proves$] Relação de dedutibilidade sintática (``a teoria prova\ldots'').
\item[$\lif$, $\liff$, $\lnot$, $\land$, $\lor$] Implicação, bicondicional, negação, conjunção, disjunção.
\item[$\lexists$, $\lforall$] Quantificadores existencial e universal.
\item[$\Th{T}$] Letra ou nome curto para uma teoria (conjunto de sentenças) de primeira ordem com símbolo primitivo~$\in$.
\end{description}

\section*{Conjuntos, relações e ordens}

\begin{description}
\item[$\emptyset$] Conjunto vazio.
\item[$\in$, $\subseteq$, $\subsetneq$] Pertence a; subconjunto; subconjunto próprio.
\item[$\cup$, $\cap$, $\setminus$] União, intersecção, diferença.
\item[$\Pow{X}$, $\wp(X)$] Conjunto potência de~$X$.
\item[$\tuple{a,b}$, $\tuple{a_1,\ldots,a_n}$] Par ordenado; $n$-upla ordenada.
\item[$\Setabs{x}{\phi(x)}$] Conjunto dos~$x$ que verificam~$\phi(x)$ (quando definível no contexto).
\item[$\equivrep{a}{R}$, $\equivclass{A}{R}$] Classe de equivalência de~$a$; conjunto das classes em~$A$.
\item[$\Id{A}$] Relação identidade sobre~$A$: $\Setabs{\tuple{x,x}}{x \in A}$.
\item[$R^n$, $R \mid S$] Potências e composição de relações (definidas no capítulo de relações).
\item[$R^+$] Fecho transitivo da relação~$R$ (união das potências $R^n$ com $n \geq 1$). \emph{Atenção:} em algumas passagens sobre ordens, $R^+$ denota $R \cup \Id{A}$ (reflexivação de uma ordem estrita); o contexto distingue.
\item[$R^*$] Fecho reflexivo transitivo: $R^+ \cup \Id{A}$ (definição padrão no capítulo de relações).
\end{description}

\section*{Funções}

\begin{description}
\item[$\dom{f}$] Domínio de~$f$.
\item[$\operatorname{Im}(f)$] Imagem de~$f$ (neste volume a imagem aparece sob esta forma).
\item[$f \restrict A$] Restrição de~$f$ ao subconjunto~$A$.
\item[$g \circ f$] Composição (aplicar primeiro~$f$, depois~$g$).
\item[$\funfromto{X}{Y}$] Conjunto de todas as funções de~$X$ para~$Y$.
\item[$\funimage{f}{A}$] Imagem do conjunto~$A$ por~$f$: $\Setabs{f(x)}{x \in A}$.
\item[$\mathrm{comp}(\sigma)$] Comprimento de uma sequência finita~$\sigma$.
\end{description}

\section*{Números habituais}

\begin{description}
\item[$\Nat$, $\Int$, $\Rat$, $\Real$] Naturais, inteiros, racionais, reais.
\end{description}

\section*{Ordinais e aritmética ordinal}

\begin{description}
\item[$\omega$] Primeiro ordinal limite (identificado com os naturais de von Neumann).
\item[$\alpha^{+}$] Sucessor ordinal de~$\alpha$.
\item[$+$, $\cdot$ entre ordinais] Adição e multiplicação \emph{ordinais} (símbolos iguais aos da aritmética usual; o contexto distingue).
\item[$\ordexpo{\alpha}{\beta}$] Exponenciação ordinal, quando aparece explicitamente.
\item[$\ordtype{A,R}$] Tipo de ordem da estrutura bem ordenada $\tuple{A,R}$.
\item[$\rlexless$] Ordem lexicográfica restrita (usada em produtos de ordinais ordenados).
\item[$\disjointsum$] Soma ordenada disjunta de ordinais (soma como ordem em cópias disjuntas).
\item[$\supstrict(X)$] Para $X$ conjunto de ordinais, o menor ordinal estritamente maior que todo o membro de~$X$ (no PDF aparece como $\mathrm{lsub}(X)$, convenção herdada do projeto original em inglês).
\end{description}

\section*{Cardinalidade e equipotência}

\begin{description}
\item[$\card{X}$] Cardinalidade (tamanho) do conjunto~$X$.
\item[$\cardeq{A}{B}$, $\cardle{A}{B}$, $\cardless{A}{B}$] Comparação de tamanhos: equipotência ($\approx$); não maior ($\preceq$); estritamente menor ($\prec$), no sentido dos capítulos sobre cardinalidade.
\item[$\cardneq{A}{B}$, $\cardnless{\alpha}{A}$] Negação de equipotência; e formulações do tipo ``não existe injeção de um ordinal em $A$'' usadas, por exemplo, no lema de Hartogs.
\item[$\prec$, $\preceq$, $\approx$] Ordem e equivalência entre cardinais (no sentido introduzido nos capítulos de cardinais).
\item[$\aleph_\alpha$, $\beth_\alpha$] Sucessões transfinitas de cardinais infinitos (\textit{alef} e \textit{beth}).
\item[$\cardfont{a}$, $\cardfont{b}$] Letras góticas para designar cardinais.
\item[$\oplus$, $\otimes$, potências cardinais] Soma, produto e exponenciação entre cardinais.
\item[$\cardsucc{\kappa}$] Sucessor cardinal, quando notado explicitamente.
\end{description}

\section*{Hierarquia cumulativa, posto e fecho}

\begin{description}
\item[$V_\alpha$] Estágio de nível~$\alpha$ na hierarquia cumulativa de conjuntos.
\item[$\mathrm{posto}(x)$] Posto do conjunto~$x$: o menor ordinal~$\alpha$ com $x \subseteq V_\alpha$ (em inglês costuma dizer-se \emph{rank}).
\item[$\mathrm{ft}(x)$] Fecho transitivo do \emph{conjunto}~$x$. A abreviatura \emph{ft} segue as iniciais de \emph{fecho transitivo}; em inglês usa-se muitas vezes \emph{trcl} (\textit{tr}ansitive \textit{cl}osure).
\item[$\closureofunder{f}{o}$] Menor conjunto que contém~$o$ e é fechado sob a função~$f$ (estruturas de Dedekind, naturais de von Neumann, etc.).
\end{description}

\section*{Objetos tipo Frege e cardinalidade Tarski--Scott}

\begin{description}
\item[$\# x\,\phi(x)$] Notação abreviada para certos objetos definidos a partir de fórmulas (número de Hume, extensões de Frege, conforme o contexto).
\item[$\mathrm{tsc}(A)$] Cardinalidade no sentido \emph{Tarski--Scott} de um conjunto~$A$.
\end{description}

\section*{Geometria e medida (paradoxos da escolha)}

\begin{description}
\item[$\onesphere$] Esfera (contexto de decomposição paradoxal).
\item[$\rotationsgroup$] Grupo de rotações usado nas construções de Vitali e afins.
\end{description}

% Glossário de termos — curso Set Theory PT (só vocabulário e noções nomeadas; ordem alfabética dentro de cada tema)

\clearpage
\phantomsection
\chapter*{Glossário}
\markboth{\MakeUppercase{Glossário}}{\MakeUppercase{Glossário}}
\addcontentsline{toc}{chapter}{Glossário}

\noindent
Este glossário reúne \textbf{termos} (palavras e expressões usadas no discurso matemático do livro). Os \textbf{símbolos} e as \textbf{listas formais} de abreviaturas teóricas estão no capítulo \emph{Notação}. As entradas agrupam-se por tema, na ordem de leitura; dentro de cada tema, vão em ordem alfabética em português.

\section*{Prelúdio e motivação histórica}

\begin{description}
\item[Continuidade] Noção de ``sem saltos'' para funções sobre o contínuo, afinada no século XIX com argumentos do tipo épsilon-delta e no rigor de Cauchy e Weierstrass; motiva a necessidade de fundamentos.
\item[Curva (Jordan)] Neste livro, mapa contínuo de um intervalo para o plano; contexto para curvas preenchedoras e exemplos geométricos ``patológicos''.
\item[Curva preenchedora de espaço] Mapa contínuo que sobrecarrega uma dimensão (por exemplo curvas de Peano ou de Hilbert), desafiando a intuição de dimensão e pressionando a noção rigorosa de função e de conjunto.
\item[Derivada] Taxa de variação instantânea, desenvolvida primeiro com ideias de infinitesimal e depois no rigor moderno.
\item[Diferenciabilidade] Propriedade de admitir derivada; contrasta com continuidade e com exemplos limite discutidos no prelúdio.
\item[Função] No prelúdio, objeto central do cálculo; mais adiante, tratada conjuntisticamente como relação funcional ou conjunto de pares ordenados.
\item[Infinitesimal] Quantidade ``infinitamente pequena'' usada heuristicamente antes do rigor moderno; parte da história que leva à formalização.
\item[Limite] Conceito que articula aproximação e continuidade; peça-chave na crise de rigor do século XIX.
\item[Mitologia da matemática] Expressão do texto para narrativas simplificadoras ou exageradas sobre descoberta e progresso; contrapõe-se ao estudo histórico sério.
\item[Patologia / patológico] Exemplos que violam intuições ingénuas (funções, curvas, cardinalidades), muitas vezes motores de novo rigor.
\item[Rigor matemático] Padrões de definição e prova que tornam os resultados auditáveis; o prelúdio explica parte da pressão histórica em direção à teoria dos conjuntos.
\item[Teoria dos conjuntos] Quadro no qual, no fim do livro, se modelam números, funções e análise; o prelúdio apresenta parte das razões históricas dessa adopção.
\end{description}

\section*{Conjuntos (teoria ingênua)}

\begin{description}
\item[Coleção extensional] Ideia de que um conjunto fica determinado pelos seus membros, não pela maneira de o descrever; base da extensionalidade.
\item[Compreensão (princípio ingénuo)] Princípio de formar o conjunto de todos os objetos que satisfazem uma condição; sob a ingenuidade total, conduz ao paradoxo de Russell.
\item[Conjunto] Coleção considerada como objeto matemático; na parte ingénua, tratada informalmente antes dos axiomas restritivos.
\item[Conjunto vazio] Conjunto sem elementos; ponto de partida de muitas construções.
\item[Conjunto potência] Conjunto de todos os subconjuntos de um conjunto dado; aparece no axioma homónimo na teoria formal.
\item[Diferença de conjuntos] Conjunto dos elementos da primeira parcela que não pertencem à segunda.
\item[Elemento / membro] Objeto que ``pertence'' a um conjunto; o texto distingue léxico conforme o original em inglês (\textit{element} / \textit{member}).
\item[Extensionalidade (ingénua)] Princípio de que conjuntos com os mesmos membros são iguais; na teoria formal torna-se o axioma de extensionalidade.
\item[Intersecção] Conjunto dos objetos que pertencem simultaneamente a todos os conjuntos de uma família; inclui intersecção de dois conjuntos e intersecção de uma família arbitrária.
\item[Par ordenado] Objeto que codifica ordem; base para relações e funções como conjuntos.
\item[Paradoxo de Russell] Contradição obtida ao considerar ``o conjunto de todos os conjuntos que não são membros de si próprios''; motiva a teoria axiomática.
\item[Parte própria] Subconjunto estrito; usado na definição de infinito de Dedekind.
\item[Produto cartesiano] Conjunto de todos os pares ordenados cujas primeiras componentes vêm de um conjunto e as segundas de outro (e generalizações).
\item[Subconjunto] Conjunto cujos membros são todos membros de outro; subconjunto próprio se, além disso, falta pelo menos um membro do conjunto maior.
\item[Tupla] Generalização de par ordenado para várias componentes; usada em aritmetização e em estruturas ordenadas.
\item[União] Conjunto dos membros dos membros de uma família; inclui união de dois conjuntos e união de uma família arbitrária.
\item[União disjunta] Construção que ``etiqueta'' cópias disjuntas de conjuntos; usada em ordinais e cardinalidade para somas ordenadas.
\end{description}

\section*{Relações e ordens (teoria ingênua)}

\begin{description}
\item[Anti-simetria] Se dois elementos se relacionam mutuamente, então são iguais; típica de ordens parciais reflexivas.
\item[Assimetria] Nunca dois elementos distintos se relacionam mutuamente nos dois sentidos; típica de ordens estritas.
\item[Classe de equivalência] Conjunto dos elementos relacionados a um dado membro por uma relação de equivalência.
\item[Conexa (relação)] Para quaisquer dois elementos distintos, um relaciona-se com o outro num dos sentidos (ordem total).
\item[Divisibilidade] Exemplo de ordem parcial sobre os inteiros; ilustra relações não lineares.
\item[Equivalência (relação)] Relação reflexiva, simétrica e transitiva; particiona o domínio em classes.
\item[Fecho reflexivo] Acrescentar à relação todos os pares na diagonal (cada elemento relacionado consigo); atenção à sobrecarga do símbolo de fecho transitivo noutros capítulos.
\item[Fecho reflexivo transitivo] Menor relação reflexiva e transitiva que contém a relação dada; no capítulo de operações aparece construída a partir do fecho transitivo e da diagonal.
\item[Fecho transitivo (relação)] Menor relação transitiva que contém a relação dada, obtida como união das potências iteradas da relação.
\item[Identidade (relação)] Relação que relaciona cada elemento apenas consigo, sobre um conjunto dado.
\item[Irreflexividade] Nenhum elemento se relaciona consigo.
\item[Ordem estrita] Relação irreflexiva e transitiva (e, em geral, assimétrica); contrapõe-se à ordem reflexiva.
\item[Ordem linear / total] Ordem parcial conexa: quaisquer dois elementos são comparáveis.
\item[Ordem parcial] Relação reflexiva, anti-simétrica e transitiva.
\item[Pré-ordem] Relação reflexiva e transitiva (não exige anti-simetria).
\item[Quociente] Conjunto de classes de equivalência sob uma relação de equivalência.
\item[Reflexividade] Todo elemento relaciona-se consigo.
\item[Relação binária] Subconjunto de um produto cartesiano; base para ordens, equivalências e funções.
\item[Relação inversa] Relação obtida invertendo a ordem dos pares.
\item[Restrição (relação)] Relação obtida considerando apenas elementos de um subconjunto do domínio.
\item[Simetria] Se um elemento relaciona-se com outro, o segundo relaciona-se com o primeiro.
\item[Transitividade] Se o primeiro relaciona-se com o segundo e o segundo com o terceiro, então o primeiro relaciona-se com o terceiro.
\end{description}

\section*{Funções (teoria ingênua)}

\begin{description}
\item[Aplicação] Termo usado quando o original distingue \textit{application} de \textit{map}; designa o valor ou o acto de aplicar uma função a um argumento.
\item[Bijeção] Função simultaneamente injetiva e sobrejetiva; admite inversa bilateral.
\item[Composição de funções] Aplicar primeiro uma função e depois outra; exige compatibilidade entre imagem da primeira e domínio da segunda.
\item[Contradomínio] Conjunto onde a função toma valores declarados; pode ser estritamente maior que a imagem.
\item[Domínio] Conjunto de argumentos para os quais a função está definida.
\item[Extensionalidade (funções)] Duas funções com o mesmo domínio e os mesmos valores ponto a ponto são iguais.
\item[Função] Relação em que cada argumento do domínio tem exactamente um valor; identificada com o seu gráfico.
\item[Gráfico (de função)] Conjunto de pares argumento-valor; identifica a função com uma relação binária especial.
\item[Imagem] Conjunto dos valores que a função toma no seu domínio.
\item[Injeção] Função que nunca envia dois argumentos distintos para o mesmo valor.
\item[Inversa] Função que reverte outra; existência global equivalente à bijeção, ou uso de inversas à esquerda ou à direita.
\item[Mapa / mapeamento] Sinónimos estáveis de \textit{map} / \textit{mapping} no original; designam função.
\item[Relação funcional] Relação em que cada elemento do domínio tem exactamente um sucessor pela relação.
\item[Restrição (função)] Função obtida limitando o domínio a um subconjunto, mantendo os mesmos valores.
\item[Sobrejeção] Função cuja imagem coincide com o contradomínio declarado.
\item[Valor] Objeto associado a um argumento pela função.
\end{description}

\section*{Cardinalidade ingênua}

\begin{description}
\item[Cardinalidade (informal)] ``Tamanho'' de um conjunto, refinado depois por equipotência e pela teoria dos cardinais.
\item[Confinito] Conjunto equipotente a um subconjunto próprio de si; na presença dos axiomas usuais, alinha-se com o infinito de Dedekind.
\item[Diagonalização] Argumento que constrói um objecto diferindo de cada linha de uma lista (Cantor, conjunto potência dos naturais, sequências binárias).
\item[Enumerável] Equipotente ao conjunto dos naturais ou a um segmento inicial finito; no texto ingénuo, também caracterizado por injeção nos naturais.
\item[Enumeração] Lista ou bijeção que organiza todos os membros de um conjunto com índices naturais.
\item[Equipotência] Relação entre dois conjuntos quando existe bijeção de um para o outro; ``mesmo tamanho'' em sentido conjuntista.
\item[Não enumerável] Conjunto que não é enumerável; exemplo paradigmático no livro: o conjunto potência dos naturais.
\item[Redução (enumerabilidade)] Provar não enumerabilidade transportando a dificuldade a partir de outro conjunto já conhecido por não enumerável.
\item[Teorema de Cantor] Para qualquer conjunto, o seu conjunto potência tem cardinalidade estritamente maior; não existe sobrejeção do conjunto para o seu conjunto potência.
\item[Teorema de Schr\"oder--Bernstein] Se cada um de dois conjuntos injeta no outro, então são equipotentes; a prova completa é adiada até à teoria dos cardinais.
\end{description}

\section*{Aritmetização e construção dos números}

\begin{description}
\item[Construção de Cauchy] Reais como classes de equivalência de sequências de Cauchy de racionais.
\item[Cortes de Dedekind] Partições dos racionais em dois conjuntos que modelam a completude da recta real.
\item[Inteiros] Construídos a partir dos naturais por equivalência em pares (diferenças formais).
\item[Racionais] Construídos como classes de equivalência de pares de inteiros (fracções formais).
\item[Reais] Construídos por cortes de Dedekind ou por Cauchy; completam a cadeia natural dos sistemas numéricos no livro.
\item[Sucessor aritmético] Operação que gera o ``próximo'' natural na estrutura deduzida dos axiomas de Peano codificados em teoria dos conjuntos.
\end{description}

\section*{Infinito, naturais e indução}

\begin{description}
\item[Álgebra de Dedekind] Estrutura com sucessor injectivo, elemento inicial fora da imagem do sucessor e princípio da indução; modelo para os naturais.
\item[Fecho indutivo] Menor conjunto que contém um gerador e é fechado sob uma operação; no tratamento de von Neumann coincide com o conjunto dos naturais.
\item[Indução matemática] Princípio que permite provar propriedades para todos os naturais a partir do caso inicial e do passo indutivo.
\item[Infinito de Dedekind] Conjunto equipotente a uma parte própria de si; caracterização do infinito antes do desenvolvimento completo da hierarquia ordinal.
\item[Números naturais] No von Neumann, o conjunto dos ordinais finitos; base para a aritmetização no livro.
\item[Recursão] Método para definir funções nos naturais (e generalizações) especificando valor inicial e passo iterativo.
\item[Sucessor (von Neumann)] Operação conjuntista que gera o próximo estágio a partir de um conjunto; usada na construção dos naturais e dos ordinais.
\end{description}

\section*{Concepção iterativa, paradoxos e filosofia}

\begin{description}
\item[Compreensão impredicativa] Definição de um objeto por referência a uma totalidade que inclui o próprio objeto; tema na crítica predicativista.
\item[Concepção iterativa] Visão de que os conjuntos são formados estágio a estágio, em contraste com a ``totalidade'' ingénua.
\item[Extensionalidade (justificação)] Motivação filosófica para identificar conjuntos pela coincidência de membros.
\item[\textit{Grundgesetze}] Obra de Frege; sistema em que a compreensão irrestrita conduz à contradição.
\item[Paradoxo de Russell (revisitado)] Retorno ao paradoxo após a narrativa axiomática; liga ingenuidade, classes e diagonalização.
\item[Predicatividade] Restrição a definições que não quantificam sobre totais ``maiores'' que o objeto definido.
\item[Urelemento] Objecto sem membros que não é o conjunto vazio; discutido como opcional face à teoria pura de conjuntos.
\end{description}

\section*{Axiomas e fragmentos de teoria}

\begin{description}
\item[Axioma da escolha] Existência de função de escolha para famílias de conjuntos não vazios; caracteriza a teoria com escolha face à sem escolha.
\item[Axioma da extensionalidade] Conjuntos com os mesmos membros são iguais.
\item[Axioma da fundamentação] Cada conjunto não vazio tem um membro disjunto dele; evita ciclos de pertença e patologias associadas.
\item[Axioma da regularidade] Formulação equivalente à fundamentação no desenvolvimento do livro (equivalência demonstrada módulo o fragmento adequado da teoria).
\item[Axioma da união] Para qualquer conjunto, existe a união dos seus membros.
\item[Axioma do conjunto potência] Para qualquer conjunto, existe o conjunto de todos os seus subconjuntos.
\item[Axioma do infinito] Afirma a existência de um conjunto indutivo adequado para fundar os naturais de von Neumann.
\item[Axioma dos pares] Para quaisquer dois objectos, existe o par (não ordenado) formado por eles.
\item[Esquema de separação] Para cada propriedade definível, o subconjunto dos membros de um conjunto dado que a verificam é conjunto; restringe a compreensão.
\item[Esquema de substituição] A imagem de um conjunto por uma relação funcional definível é conjunto; distingue a força de Zermelo--Fraenkel da de Zermelo puro.
\item[Teoria de Zermelo--Fraenkel] Teoria de primeira ordem com predicado de pertença, axiomas de Zermelo e substituição; variantes sem escolha ou com escolha aparecem no texto.
\item[Teoria de Zermelo] Fragmento com os axiomas de Zermelo (sem substituição, conforme o enquadramento do livro).
\item[Teoria dos Níveis] Teoria mais fraca baseada na ideia de ``níveis''; relacionada com Zermelo--Fraenkel e com extensões discutidas na parte de substituição.
\item[Extensão Zr] Teoria obtida acrescentando axiomas fortes à Teoria dos Níveis; no texto, todo conjunto tem posto e a teoria supera Zermelo puro.
\end{description}

\section*{Ordinais}

\begin{description}
\item[Boa ordenação] Ordem linear em que todo subconjunto não vazio tem mínimo; em certos contextos, equivalente a princípios de escolha.
\item[Indução transfinita] Generalização da indução para ordinais arbitrários; ferramenta central em provas sobre a hierarquia e sobre postos.
\item[Limite (ordinal)] Ordinal não zero que não é sucessor; acumula todos os ordinais menores.
\item[Ordinal] No sentido de von Neumann: conjunto transitivo bem ordenado pela pertença.
\item[Ordinal inicial] Menor ordinal de uma dada cardinalidade infinita; liga cardinais a ordinais escolhidos canonicamente.
\item[Ordem bem ordenada] Estrutura ordenada isomorfa a um ordinal, ou, equivalentemente, estritamente bem fundamentada na ordem.
\item[Representação ordinal] Todo bem ordenamento é isomorfo a um único ordinal.
\item[Segmento inicial] Subconjunto fechado para predecessores na ordem; usado em isomorfismos de ordens.
\item[Sucessor ordinal] Próximo estágio após um ordinal dado na construção de von Neumann.
\item[Transitivo (conjunto)] Se um membro pertence a um membro do conjunto, então pertence ao próprio conjunto; propriedade dos ordinais e dos estágios da hierarquia.
\end{description}

\section*{Hierarquia cumulativa, posto e fundação}

\begin{description}
\item[Estágio da hierarquia] Nível indexado por um ordinal na hierarquia cumulativa; conjunto de tudo o que aparece ``antes'' desse ordinal.
\item[Fecho transitivo (conjunto)] Menor conjunto transitivo que contém um conjunto dado; no volume, o texto adopta uma abreviatura consagrada na lista de notação.
\item[Hierarquia cumulativa] Família de estágios indexada pelos ordinais, acumulando conjuntos estágio a estágio; modelo formal da concepção iterativa.
\item[Posto] Menor ordinal $\alpha$ tal que o conjunto dado é subconjunto do estágio $\alpha$ da hierarquia; mede profundidade na cumulativa.
\item[Recursão na hierarquia] Definições e provas por indução nos ordinais ou nos estágios da hierarquia.
\end{description}

\section*{Substituição e limitação de tamanho}

\begin{description}
\item[Limitação de tamanho] Princípio que impede a formação de ``conjuntos demasiado grandes'' sem critério; contrapõe-se à compreensão irrestrita.
\item[Reflexão (metatema)] Tópicos sobre o que a linguagem de primeira ordem consegue exprimir sobre a hierarquia vista como totalidade.
\end{description}

\section*{Aritmética ordinal}

\begin{description}
\item[Adição ordinal] Operação que ``cola'' duas boas ordenações; em geral não comutativa.
\item[Exponenciação ordinal] Potenciação transfinita baseada na multiplicação ordinal; distinta da exponenciação cardinal.
\item[Multiplicação ordinal] Repetição ordenada de uma boa ordenação; liga-se a produtos lexicográficos.
\item[Ordem lexicográfica] Ordem em produtos que compara primeiro numa coordenada e depois noutra; usada na aritmética ordinal.
\end{description}

\section*{Cardinais e comparabilidade}

\begin{description}
\item[Cardinal] No livro, sobretudo como ordinal inicial ou como objeto que codifica o ``tamanho'' de um conjunto; as sucessões alef e beth dos infinitos aparecem na parte final.
\item[Cardinal inicial] Ver \emph{Ordinal inicial}; identifica um cardinal infinito com o menor ordinal da sua cardinalidade.
\item[Comparabilidade de cardinais] Relações de não maior, estritamente menor e equipotência entre tamanhos; sem Axioma da Escolha, a comparabilidade falha em geral.
\item[Hartogs (ordinal de)] Menor ordinal que não se injeta num conjunto dado; ferramenta para comparar tamanhos sem escolha global.
\item[Tipo de ordem] Ordinal associado a um bem ordenamento, único a menos de isomorfismo.
\end{description}

\section*{Aritmética cardinal e hipóteses do contínuo}

\begin{description}
\item[Aritmética cardinal] Operações de soma, produto e potência entre cardinais; leis distintas da aritmética ordinal.
\item[Cofinalidade] Menor ordinal que serve de domínio para uma função com imagem ilimitada num cardinal dado; usada em limitações de cardinalidade.
\item[Continuo] Cardinalidade da recta real; no desenvolvimento clássico, identifica-se com a potência do conjunto dos naturais.
\item[Hipótese do Contínuo] Não existe cardinal estritamente entre o dos naturais e o do continuo; independente da teoria usual com escolha.
\item[Hipótese Generalizada do Contínuo] Identificação, para todo ordinal, das duas sucessões canónicas dos infinitos que o livro nota com alef e beth; implica a Hipótese do Contínuo.
\item[Independência da hipótese do contínuo] Resultado de que nem a hipótese nem a sua negação se demonstram só com os axiomas habituais; o livro resume o panorama (Gödel, Cohen e posteriores).
\item[Ponto fixo de alef ou de beth] Cardinal que coincide com o termo indexado pela própria cardinal na sucessão alef ou beth; existência por argumentos de limite.
\item[Sucessor cardinal] Menor cardinal estritamente maior que um cardinal dado.
\end{description}

\section*{Escolha e princípios fracos}

\begin{description}
\item[Axioma da boa ordenação] Todo conjunto pode ser bem ordenado por alguma relação; no livro figura na enumeração da teoria com escolha e é equivalente ao Axioma da Escolha.
\item[Conjunto não mensurável] Subconjunto da recta real sem medida de Lebesgue compatível com a invariância por translação; construção associada a escolhas (Vitali).
\item[Decomposição paradoxal] Construções que decompõem sólidos e recompõem por isometrias, possíveis com o Axioma da Escolha (contexto Banach--Tarski).
\item[Escolha enumerável] Princípio que afirma escolha apenas para famílias enumeráveis de conjuntos não vazios; capítulo dedicado.
\item[Escolha (noção)] Uso frequente do texto para designar o Axioma da Escolha quando se trabalha na teoria com escolha.
\item[Função de escolha] Associa a cada membro de uma família de conjuntos não vazios um elemento desse membro; a existência global para toda família é o Axioma da Escolha.
\item[Teorema da boa ordenação] Todo conjunto admite uma boa ordenação; equivalente ao Axioma da Escolha no tratamento clássico.
\end{description}

\section*{Cardinalidade no sentido de Tarski e Scott}

\begin{description}
\item[Cardinalidade Tarski--Scott] Noção de ``tamanho'' que dispensa o teorema da boa ordenação; permite falar de cardinalidade sem assumir o Axioma da Escolha.
\item[Número de Hume] Abstracção de cardinalidade por equinumerosidade, no enquadramento filosófico do texto.
\item[Extensão de Frege] Objeto abstrato associado a uma condição; contexto histórico e lógico ligado a paradoxos e a alternativas conjuntistas.
\end{description}


\olimport*[]{open-logic-about}